\label{sec:discussion}

\subsection{Party-Specific Analysis}

The safe-seat comparison reveals an asymmetric pattern across the three states. In NC
and TX---states where Republicans controlled redistricting in 2021--2022---the enacted
maps produce significantly more safe-Republican seats than the algorithmic baseline,
while the safe-Democratic seat count is comparable or lower. In WI, where the
Republican-controlled legislature's map was partially modified by court intervention,
the asymmetry is smaller.

This asymmetry reflects the mechanism of partisan gerrymandering: the party in control
of redistricting packs opposition voters into very safe seats (wasting their votes in
blowout victories) while cracking them across safe-majority seats (denying them
representation). The resulting map has more safe seats overall, but the additional safe
seats are predominantly for the party in control.

Table~\ref{tab:party-decomp} reports the party-specific safe-seat counts for all three
states, highlighting the asymmetric nature of the safe-seat creation.

\begin{table}[ht]
\centering
\begin{threeparttable}
\caption{Party-specific safe-seat decomposition.}
\label{tab:party-decomp}
\begin{tabular}{l l r r r}
\toprule
State & Plan & Safe-D seats & Safe-R seats & Difference \\
\midrule
NC ($k=14$) & \textsc{bisect}$^\dagger$ & 2 & 6 & $-$4 (R advantage) \\
            & Enacted 2022              & 3 & 9 & $-$6 (R advantage) \\
            & Additional (enacted--alg) & +1 & +3 & \\
\addlinespace
WI ($k=8$)  & \textsc{bisect}$^\dagger$ & 2 & 1 & +1 (D advantage) \\
            & Enacted 2022              & 2 & 2 & 0 \\
            & Additional (enacted--alg) & 0 & +2 & \\
\addlinespace
TX ($k=38$) & \textsc{bisect}$^\dagger$ & 8 & 14 & $-$6 (R advantage) \\
            & Enacted 2022              & 7 & 19 & $-$12 (R advantage) \\
            & Additional (enacted--alg) & $-$1 & +5 & \\
\bottomrule
\end{tabular}
\begin{tablenotes}
\small
\item[$\dagger$] Single-run result (seed\,$=42$).
\end{tablenotes}
\end{threeparttable}
\end{table}

\subsection{Geographic Determinism of Democratic Safe Seats}

A consistent finding across all three states is that safe-Democratic seats are largely
preserved by geographic compactness: \textsc{bisect} maps produce safe-D seat counts
comparable to enacted maps (or even slightly higher in TX) because urban Democratic
voter concentrations are sufficiently dense that any compact district centred on them
will produce a safe-D seat.

This geographic determinism of Democratic safe seats has an important implication for
the free-elections doctrine analysis discussed below: the urban-core safe-D seats in
\textsc{bisect} maps are not a product of packing strategy but of geographic sorting.
Their presence in both algorithmic and enacted maps confirms that they cannot be
eliminated without violating compactness, creating prison gerrymandering concerns, or
fragmenting cities.

The additional safe seats in enacted maps are predominantly safe-Republican seats in
suburban and exurban areas, where the geographic partisan mix would, under compact
redistricting, produce competitive or lean districts. Converting these lean-R or
competitive tracts into safe-R districts requires extending district boundaries to
incorporate more rural Republican precincts, creating less compact shapes.

\subsection{Legal Implications: Free Elections Doctrine}

The free elections doctrine analysis builds on the \textsc{bisect} comparison to
establish that the enacted maps' safe-seat counts substantially exceed the geographic
minimum. Compact redistricting (as operationalised by the algorithmic map) produces
$X$ safe seats; the enacted map produces $X + \Delta$ safe seats, where $\Delta > 0$
measures the enacted map's safe-seat premium. This premium represents the additional
electoral insulation that incumbents receive from deliberate district boundary choices.

State supreme courts evaluating free-elections claims---following the precedents in
\textit{League of Women Voters v.\ Commonwealth}\footnote{\textit{League of Women
Voters v.\ Commonwealth}, 178 A.3d 737 (Pa.\ 2018).} and \textit{Harper v.\ Hall}
\footnote{\textit{Harper v.\ Hall}, 868 S.E.2d 499 (N.C.\ 2022).}---have focused
precisely on this question: does the enacted map suppress competitive elections beyond
what geographic necessity requires? The \textsc{bisect} comparison provides the
counterfactual that courts need to answer this question: a geographically compact plan
that produces more competitive districts than the enacted map demonstrates that the
enacted map's low competitive-seat count is a choice, not a necessity.

\subsection{Compactness as an Indirect Cause}

The mechanism by which \textsc{bisect} produces fewer safe seats than enacted maps is
indirect and important: compactness optimisation, not any explicit competitive-district
target, drives the result. The algorithm does not know about partisan lean; it draws
compact districts without regard to whether those districts are safe or competitive.
But compact districts in suburban transition zones---where Democratic and Republican
voters are geographically interspersed---produce competitive seat shares that deliberate
packing and cracking strategies avoid.

This indirect effect has a legal implication: \textsc{bisect}'s more competitive maps
are not a consequence of the algorithm pursuing competitiveness as a goal (which would
be a permissible but non-required redistricting criterion under state law). They are
a consequence of the algorithm pursuing compactness, which is both a traditional
redistricting criterion and constitutionally required in many states. The competitiveness
of \textsc{bisect} maps is a side effect of geographic optimisation, not a target.
